\documentclass[11pt]{amsart}
\usepackage{amsmath,amssymb}
\setlength{\emergencystretch}{3em}
\tolerance=600
\usepackage{tikz-cd}
\usepackage[hidelinks]{hyperref}

\newtheorem{thm}{Theorem}[section]
\newtheorem{lem}[thm]{Lemma}
\newtheorem{cor}[thm]{Corollary}
\newtheorem{prop}[thm]{Proposition}
\theoremstyle{definition}
\newtheorem{defi}[thm]{Definition}
\newtheorem{rem}[thm]{Remark}
\newtheorem{ex}[thm]{Example}
\newtheorem{que}[thm]{Question}

\newcommand{\Same}{\mathrm{Same}}
\newcommand{\Sub}{\mathrm{Sub}}
\newcommand{\Mor}{\mathrm{Mor}}
\newcommand{\Iso}{\mathrm{Iso}}
\newcommand{\Con}{\mathrm{Con}}
\newcommand{\Cg}{\mathrm{Cg}}
\newcommand{\Hom}{\mathrm{Hom}}
\newcommand{\Presh}{\mathrm{Presh}}
\newcommand{\NExt}{\mathrm{NExt}}
\newcommand{\Ext}{\mathrm{Ext}}
\newcommand{\sV}{\mathcal V}
\newcommand{\sW}{\mathcal W}
\newcommand{\sR}{\mathcal R}
\newcommand{\sS}{\mathcal S}
\newcommand{\lT}{\lambda\mathcal T}
\newcommand{\var}{\mathrm{var}}

\begin{document}

\title[Subvarieties and logics in the sameness lattice]{Lattices of subvarieties and of logics as sublattices of the lattice of two-out-of-three classes}
\author{Ryota Shibusawa}
\thanks{Corresponding author. E-mail: \texttt{r-shibusawa@daiichi-koudai.ac.jp}}
\address{Daiichi Institute of Technology}
\email{r-shibusawa@daiichi-koudai.ac.jp}
\keywords{variety, subvariety lattice, verbal congruence, reflective subcategory,
  two-out-of-three, congruence distributive, congruence permutable, Jónsson's lemma,
  lattice of modal logics, substructural logic, lambda theories}
\subjclass[2020]{08B15 (primary), 08B10, 18A40, 03B45, 03B47, 03B40 (secondary)}

\begin{abstract}
For a category $C$ let $\Same(C)$ be the complete lattice of classes of
morphisms that contain the identities and are closed under composition and
two-out-of-three. Every reflective subcategory $\sR$ contributes the element
$W_{\sR}$ of maps inverted by the reflector. We study this assignment on the
lattice $\Sub(\sV)$ of subvarieties of a variety $\sV$ of algebras. The meet of
two subvarieties always goes to the join of the two classes,
$W_{V\wedge V'}=W_V\vee W_{V'}$, because the reflector into $V\wedge V'$ is the
composite of the two reflectors; for Grothendieck topologies the corresponding
statement requires towers of plus constructions and can fail. The join of two
subvarieties goes to the meet of the classes, $W_{V\vee V'}=W_V\cap W_{V'}$,
whenever $\sV$ is congruence distributive and congruence permutable, by
Jónsson's lemma and a Chinese remainder argument; so for arithmetical
varieties $V\mapsto W_V$ is a lattice embedding of $\Sub(\sV)^{\mathrm{op}}$
into $\Same(\sV)$. Neither hypothesis can be dropped: the modular group of order
$16$ gives a counterexample in groups, and the three-element chain inside the
four-element Boolean lattice gives one in bounded lattices with a constant.
Through algebraizability the lattices of normal modal logics, of intermediate
logics, of substructural logics over $\mathbf{FL}$ and of fuzzy logics are
therefore sublattices of the sameness lattices of their algebras; the lattice
of $\lambda$-theories embeds as a join-semilattice, and whether it embeds as a
lattice is left open.
\end{abstract}

\maketitle

\section{Introduction}

A \emph{sameness} on a category $C$ is a class of morphisms containing the
identities and closed under composition and two-out-of-three; the samenesses
form a complete lattice $\Same(C)$ under inclusion, with meets given by
intersection and joins by closure \cite{Same}. Every functor $F$ out of $C$
gives a sameness $W_F=\{f: Ff\text{ invertible}\}$, and in earlier work the
elements $W_{a_J}$ contributed by the sheafification functors $a_J$ of the
Grothendieck topologies $J$ on a small category were studied: $J\mapsto W_{a_J}$
preserves all meets, preserves binary joins for every finite category, and
fails to preserve them on the poset of natural numbers
\cite{Topjoin,Finitejoin}. The positive results are proved by showing that
alternating towers of plus constructions, or of separated quotients, connect a
presheaf to its sheafification for the join topology inside the join of the two
classes.

This paper carries out the same programme for the other classical lattice of
reflective subcategories: the lattice $\Sub(\sV)$ of subvarieties of a variety
$\sV$ of algebras. Each subvariety $V$ is reflective, with reflector
$a_V(A)=A/\theta_V(A)$ for the verbal congruence $\theta_V(A)$, and we write
$W_V:=W_{a_V}=\{f:a_Vf\text{ is an isomorphism}\}$. The two lattice operations
behave differently.

\begin{thm}\label{thm:A}
For every variety $\sV$ and all subvarieties $V,V'$,
$W_{V\wedge V'}=W_V\vee W_{V'}$ in $\Same(\sV)$.
\end{thm}

\begin{thm}\label{thm:B}
For subvarieties $V,V'$ of any variety, $W_{V\vee V'}=W_V\cap W_{V'}$ if and only if
\begin{enumerate}
\item[(i)] every algebra of $V\vee V'$ is a subdirect product of its $V$-reflection and
its $V'$-reflection, and
\item[(ii)] no algebra of $V\vee V'$ has a proper subalgebra that is both $V$-dense and
$V'$-dense, where a subalgebra is $V$-dense if its inclusion is inverted by the
$V$-reflector.
\end{enumerate}
Condition (ii) holds whenever the verbal congruences $\theta_V(A)$ and $\theta_{V'}(A)$
permute for every $A\in V\vee V'$; together with (i) this says that $A$ is the pullback
of its $V$- and $V'$-reflections over its $V\wedge V'$-reflection.
\end{thm}

\begin{cor}\label{cor:arith}
If $\sV$ is congruence distributive and congruence permutable, then
$W_{V\vee V'}=W_V\cap W_{V'}$ for all subvarieties $V,V'$, and
$V\mapsto W_V$ is a lattice embedding $\Sub(\sV)^{\mathrm{op}}\hookrightarrow\Same(\sV)$.
Neither hypothesis can be omitted: the equality fails for the subvarieties
$\var(\mathbb Z_8)$ and $\var(D_8)$ of groups (Example~\ref{ex:groups}), and
for the subvarieties $c=0$ and $c=1$ of bounded lattices with a constant $c$
(Example~\ref{ex:lattices}); permutability, on the other hand, is not necessary
(Example~\ref{ex:unary}).
\end{cor}

Theorem~\ref{thm:A} is elementary: because a subvariety is closed under
quotients, the reflector into $V\wedge V'$ is the composite $a_{V'}a_V$, so the
unit of $V\wedge V'$ is a composite of a $V$-unit and a $V'$-unit and the join
theorem of \cite{Postnikov} applies. The contrast with sheaf theory is the
point: sheafification for $J'$ does not preserve $J$-sheaves, the reflector
into $J\vee J'$ is not a composite of the two, and the towers of
\cite{Topjoin,Finitejoin} are the substitute, which exists for finite
categories and does not exist on $\mathbb N$. Theorem~\ref{thm:B} is where the
algebra enters. A map $g$ of $V\vee V'$ lies in $W_{V\vee V'}$ only if it is an
isomorphism, while $g\in W_V\cap W_{V'}$ only sees the two reflections of $g$;
so the question is whether the pair of reflectors $(a_V,a_{V'})$ reflects
isomorphisms on $V\vee V'$, i.e.\ whether the ``fracture'' functor
$V\vee V'\to V\times_{V\wedge V'}V'$ of Figure~\ref{fig:fracture} is
conservative. Condition (i) is injectivity, (ii) surjectivity. In a
congruence distributive variety (i) holds for all pairs by Jónsson's lemma
\cite{Jonsson}, since the subdirectly irreducible algebras of $V\vee V'$ lie in
$V\cup V'$; in a congruence permutable variety (ii) holds by a Chinese remainder
argument. In an arithmetical variety every algebra of $V\vee V'$ is therefore
the pullback of its $V$- and $V'$-reflections over its $V\wedge V'$-reflection,
the algebraic analogue of a fracture square.

\subsection*{What is known.}
The reflector of a subvariety, verbal congruences, Jónsson's lemma and the
Mal'cev condition for permutability are standard \cite{BS}. That a join of two
subvarieties of a congruence permutable variety is the direct product of the
two when they are disjoint is due to Jónsson and Tsinakis; joins that are
subdirect products are studied by Kowalski and Paoli \cite{KP}. The lattice of
samenesses, the join theorem and the results for Grothendieck topologies are
from \cite{Same,Postnikov,Topjoin,Finitejoin}.

\subsection*{What is new.}
Theorems~\ref{thm:A} and~\ref{thm:B}, the three examples, and the
consequences for logic (Section~\ref{sec:logic}): since the lattice of
axiomatic extensions of an algebraizable logic is dually isomorphic to the
subvariety lattice of its equivalent algebraic semantics \cite{BP}, the lattices
of normal modal logics, of intermediate logics, of substructural logics over
$\mathbf{FL}$ and of the usual fuzzy logics are sublattices of the sameness
lattices of modal algebras, Heyting algebras, $\mathbf{FL}$-algebras and
their subvarieties (Corollary~\ref{cor:logics}); the lattice of
$\lambda$-theories, which is isomorphic to the lattice of equational theories
of lambda abstraction algebras \cite{Salibra00}, embeds as a join-semilattice
(Corollary~\ref{cor:lambda}), and the question whether it embeds as a lattice is
left open, as lambda abstraction algebras satisfy no non-trivial congruence
lattice identity \cite{LS04}.

\section{Preliminaries}\label{sec:prelim}

\subsection{Samenesses.}
For a category $C$, $\Same(C)$ is the set of samenesses ordered by inclusion;
it is a complete lattice, $U\wedge U'=U\cap U'$, and $U\vee U'$ is the smallest
sameness containing $U\cup U'$ \cite{Same}. For a functor $P\colon C\to C$ put
$W_P=\{f:Pf\text{ invertible}\}$; it is a sameness containing all
isomorphisms.

\begin{thm}[join theorem, {\cite[Thm.~4.1]{Postnikov}}]\label{thm:join}
Let $U,V$ be samenesses on $C$ whose join contains the isomorphisms, and let
$P\colon C\to C$ be a functor with a natural transformation
$\pi\colon\mathrm{id}\Rightarrow P$ such that $\pi_X\in U\vee V$ for all $X$
and $P$ inverts $U$ and $V$. Then $U\vee V=W_P$.
\end{thm}
\begin{proof}
$W_P\supseteq U\vee V$ since $W_P$ is a sameness containing $U$ and $V$. If
$Pf$ is invertible then $Pf\in U\vee V$, so the naturality square
$\pi_Yf=Pf\,\pi_X$ has three sides in $U\vee V$, hence the fourth.
\end{proof}

\subsection{Subvarieties and verbal congruences.}
Let $\sV$ be a variety of algebras of a fixed signature and $\Sub(\sV)$ its
lattice of subvarieties; $V\wedge V'=V\cap V'$ and $V\vee V'$ is the variety
generated by $V\cup V'$. For $A\in\sV$ and $V\in\Sub(\sV)$ the \emph{verbal
congruence} $\theta_V(A)$ is the smallest congruence of $A$ whose quotient lies
in $V$; it is the congruence generated by the pairs $(s(\bar a),t(\bar a))$ for
the identities $s=t$ of $V$ and tuples $\bar a$ of $A$, and a homomorphism
$f\colon A\to B$ satisfies $f(\theta_V(A))\subseteq\theta_V(B)$. Thus
$a_V(A)=A/\theta_V(A)$ is a reflector of $\sV$ onto $V$ with unit the quotient
map $u_A\colon A\to a_VA$, and we put
\[
W_V:=\{f\in\Mor(\sV):a_Vf\text{ is an isomorphism}\}\in\Same(\sV).
\]
If $V\subseteq V'$ then $a_V=a_Va_{V'}$ and $W_{V'}\subseteq W_V$. The map
$V\mapsto W_V$ is an order-embedding $\Sub(\sV)^{\mathrm{op}}\hookrightarrow\Same(\sV)$:
if $V\not\subseteq V'$, pick $A\in V\setminus V'$; its $V'$-unit is in
$W_{V'}$, not an isomorphism, and $a_V$ sends it to itself (as $a_VA=A$ and
$a_V(A/\theta_{V'}(A))=A/\theta_{V'}(A)$ because both lie in $V$, which is closed under quotients), so it is not in
$W_V$. We write $\Mor(\sV)$ for the class of homomorphisms between algebras of
$\sV$.

\begin{lem}\label{lem:composite}
For $V,V'\in\Sub(\sV)$ and $A\in\sV$ the composite $A\to A/\theta_V(A)\to
(A/\theta_V(A))/\theta_{V'}(A/\theta_V(A))$ is the $V\wedge V'$-unit of $A$;
that is, $a_{V\wedge V'}=a_{V'}a_V$ and
$\theta_{V\wedge V'}(A)=\theta_V(A)\vee\theta_{V'}(A)$.
\end{lem}
\begin{proof}
The target lies in $V'$ by construction and in $V$ because $V$ is closed under
quotients. Any homomorphism $A\to C$ with $C\in V\cap V'$ factors through
$A/\theta_V(A)$ since $C\in V$, and the factor through
$(A/\theta_V(A))/\theta_{V'}$ since $C\in V'$. By the correspondence theorem the
kernel of the composite is $\theta_V(A)\vee\theta_{V'}(A)$.
\end{proof}

\section{Meets of subvarieties}\label{sec:meets}

\begin{proof}[Proof of Theorem~\ref{thm:A}]
Apply Theorem~\ref{thm:join} with $U=W_V$, $V=W_{V'}$ (the join contains the
isomorphisms because each does), $P=a_{V\wedge V'}$ and $\pi$ the unit. By
Lemma~\ref{lem:composite}, $\pi_A$ is the composite of the $V$-unit of $A$,
which lies in $W_V$ because $a_V$ of a unit is an isomorphism, and the
$V'$-unit of $A/\theta_V(A)$, which lies in $W_{V'}$; so $\pi_A\in W_V\vee
W_{V'}$. Since $a_{V\wedge V'}=a_{V\wedge V'}a_V=a_{V\wedge V'}a_{V'}$, the
functor $a_{V\wedge V'}$ inverts $W_V$ and $W_{V'}$. Hence
$W_V\vee W_{V'}=W_{a_{V\wedge V'}}=W_{V\wedge V'}$.
\end{proof}

\begin{figure}[ht]
\[
\begin{tikzcd}[column sep=large]
A\arrow[r,"u_A\in W_V",two heads]\arrow[rr,bend right=25,"\pi_A"'] & A/\theta_V(A)\arrow[r,"u'\in W_{V'}",two heads] & A/(\theta_V(A)\vee\theta_{V'}(A))=a_{V\wedge V'}A
\end{tikzcd}
\]
\caption{The unit of $V\wedge V'$ is a composite of two units (Lemma~\ref{lem:composite}); this
single step replaces the towers needed for Grothendieck topologies.}
\end{figure}

The argument uses only that each subvariety is closed under the reflections
into the other; it holds verbatim for reflective subcategories $\sR,\sR'$ of
any category with $a_{\sR'}(\sR)\subseteq\sR$, where then
$a_{\sR\cap\sR'}=a_{\sR'}a_{\sR}$ and $W_{\sR}\vee W_{\sR'}=W_{\sR\cap\sR'}$.
Full subcategories of sheaves do not have this property: for two Grothendieck
topologies $J,J'$ the $J'$-sheafification of a $J$-sheaf need not be a
$J$-sheaf, $a_{J\vee J'}$ is not the composite $a_{J'}a_J$, and the join
$W_{a_J}\vee W_{a_{J'}}$ is the class $W_{a_{J\vee J'}}$ only when a tower of
alternating constructions reaches the join sheafification, which it does for
finite categories \cite{Finitejoin} and does not on the poset $\mathbb N$
\cite[Thm.~7.1]{Topjoin}.

\section{Joins of subvarieties}\label{sec:joins}

The inclusion $W_{V\vee V'}\subseteq W_V\cap W_{V'}$ always holds, since
$a_V=a_Va_{V\vee V'}$ and $a_{V'}=a_{V'}a_{V\vee V'}$. For the converse we
reduce to maps inside $V\vee V'$.

\begin{lem}\label{lem:reduce}
$W_{V\vee V'}=W_V\cap W_{V'}$ holds if and only if every homomorphism
$g\colon A\to B$ with $A,B\in V\vee V'$ for which $a_Vg$ and $a_{V'}g$ are
isomorphisms is itself an isomorphism.
\end{lem}
\begin{proof}
If $f\in W_V\cap W_{V'}$, then $g:=a_{V\vee V'}f$ is a map in $V\vee V'$ with
$a_Vg=a_Vf$ and $a_{V'}g=a_{V'}f$ isomorphisms; if the condition holds, $g$ is an
isomorphism and $f\in W_{V\vee V'}$. Conversely, a map $g$ in $V\vee V'$ has
$a_{V\vee V'}g=g$, so $g\in W_{V\vee V'}$ iff $g$ is an isomorphism.
\end{proof}

\begin{defi}
A subalgebra $C\subseteq B$ is \emph{$V$-dense} if the inclusion lies in $W_V$,
i.e.\ if $C$ meets every class of $\theta_V(B)$ and $\theta_V(B)\cap C^2=\theta_V(C)$.
\end{defi}

\begin{proof}[Proof of Theorem~\ref{thm:B}]
Suppose (i) and (ii) and let $g\colon A\to B$ be a map in $V\vee V'$ with $a_Vg$,
$a_{V'}g$ isomorphisms; by Lemma~\ref{lem:reduce} we must show that $g$ is
bijective. If $g(x)=g(y)$ then $x\,\theta_V(A)\,y$ and $x\,\theta_{V'}(A)\,y$
because $a_Vg$ and $a_{V'}g$ are injective, so $x=y$ by (i): $g$ is injective.
Its image $C:=g(A)$ is a subalgebra of $B$ with $C\cong A$, and the inclusion
$C\hookrightarrow B$ is inverted by $a_V$ and $a_{V'}$ (it is $g$ composed with
an isomorphism); so $C$ is $V$-dense and $V'$-dense and $C=B$ by (ii).

Conversely suppose $W_{V\vee V'}=W_V\cap W_{V'}$. If (i) fails for some
$A\in V\vee V'$, let $\theta:=\theta_V(A)\cap\theta_{V'}(A)\ne\Delta$ and
$g\colon A\to B:=A/\theta$ the quotient map; $B\in V\vee V'$ as a subdirect
product of $A/\theta_V(A)\in V$ and $A/\theta_{V'}(A)\in V'$. By the
correspondence theorem $\theta_V(B)=(\theta_V(A)\vee\theta)/\theta
=\theta_V(A)/\theta$, so $a_Vg$ is the identity of $A/\theta_V(A)$, and likewise
for $V'$; thus $g\in W_V\cap W_{V'}$, but $g$ is not injective, contradicting
Lemma~\ref{lem:reduce}. If (ii) fails, a $V$-dense and $V'$-dense proper
subalgebra $C\subsetneq B$ with $B\in V\vee V'$ gives an inclusion in
$W_V\cap W_{V'}$ which is not an isomorphism, again contradicting
Lemma~\ref{lem:reduce}.

Finally let the verbal congruences permute on every algebra of $V\vee V'$ and
let $C\subseteq B$ be $V$-dense and $V'$-dense, $B\in V\vee V'$; we show $C=B$.
Write $\theta_V,\theta_{V'}$ for the verbal congruences of $B$. Given $b\in B$,
density gives $c_1,c_2\in C$ with $c_1\,\theta_V\,b$ and $c_2\,\theta_{V'}\,b$.
Then $c_1$ and $c_2$ are related by $\theta_V\vee\theta_{V'}=\theta_{V\wedge V'}$
(Lemma~\ref{lem:composite}), hence, as $a_{V\wedge V'}=a_{V'}a_V$ inverts the
inclusion, by $\theta_{V\wedge V'}(C)=\theta_V(C)\vee\theta_{V'}(C)$ inside $C$;
by permutability in $C$ (also an algebra of $V\vee V'$) there is $c\in C$ with
$c_1\,\theta_V(C)\,c\,\theta_{V'}(C)\,c_2$. Then $c\,\theta_V\,c_1\,\theta_V\,b$
and $c\,\theta_{V'}\,c_2\,\theta_{V'}\,b$, so $c\,(\theta_V\cap\theta_{V'})\,b$
and $c=b$ by (i) for $B$.
\end{proof}

\begin{figure}[ht]
\[
\begin{tikzcd}
A\arrow[r,two heads]\arrow[d,two heads]\arrow[dr,phantom,"\lrcorner",very near start] & A/\theta_V(A)\arrow[d,two heads]\\
A/\theta_{V'}(A)\arrow[r,two heads] & A/(\theta_V(A)\vee\theta_{V'}(A))
\end{tikzcd}
\]
\caption{The fracture square of $A\in V\vee V'$. Condition (i) of Theorem~\ref{thm:B}
is injectivity of the comparison map
$A\to A/\theta_V\times_{A/(\theta_V\vee\theta_{V'})}A/\theta_{V'}$, and
permutability of $\theta_V(A)$, $\theta_{V'}(A)$ is its surjectivity; the
square is a pullback for every $A$ in an arithmetical variety.}
\label{fig:fracture}
\end{figure}

\begin{proof}[Proof of Corollary~\ref{cor:arith}]
Let $\sV$ be congruence distributive and congruence permutable. The verbal
congruences permute, so (ii) holds by Theorem~\ref{thm:B}. For (i), let
$A\in V\vee V'$. By Jónsson's lemma \cite{Jonsson} every subdirectly
irreducible algebra of $V\vee V'$ lies in $HSP_u(V\cup V')$; an ultraproduct of
algebras each in $V$ or in $V'$ lies in $V$ or in $V'$ (one of the two index
sets belongs to the ultrafilter), and $V$, $V'$ are closed under $H$ and $S$, so
the subdirectly irreducibles of $V\vee V'$ lie in $V\cup V'$. Writing $A$ as a
subdirect product of subdirectly irreducible quotients $A/\kappa_i$, each
$\kappa_i$ contains $\theta_V(A)$ or $\theta_{V'}(A)$, hence contains
$\theta_V(A)\cap\theta_{V'}(A)$; as $\bigcap_i\kappa_i$ is the identity, so is
$\theta_V(A)\cap\theta_{V'}(A)$, which is (i). Theorem~\ref{thm:B} gives
$W_{V\vee V'}=W_V\cap W_{V'}$, and with Theorem~\ref{thm:A} and the
order-embedding of Section~\ref{sec:prelim} the map $V\mapsto W_V$ is a lattice
embedding. The examples below show that the hypotheses cannot be dropped and
that permutability is not necessary.
\end{proof}

\begin{ex}[permutable, not distributive]\label{ex:groups}
In the variety of groups let $V=\var(\mathbb Z_8)$, the abelian groups of
exponent dividing $8$, and $V'=\var(D_8)$, the variety generated by the
dihedral group of order $8$. Let $M_{16}=\langle a,b\mid a^8=b^2=1,\
bab=a^5\rangle$ be the modular group of order $16$. It lies in $V\vee V'$: in
$\mathbb Z_8\times D_8$, with $r$ a rotation of order $4$ and $s$ a reflection,
the subgroup $K$ generated by $x=(1,r)$ and $y=(0,s)$ contains
$x\cdot yxy=(1,r)(1,r^{-1})=(2,1)$, hence $x^2(2,1)^{-1}=(0,r^2)$ and
$n:=x^4(0,r^2)=(4,r^2)$, a central element of order $2$. In $K/\langle n\rangle$
the image of $x$ has order $8$ (as $x^4=(4,1)\notin\langle n\rangle$), the image
of $y$ has order $2$, and $yxy=(1,r^{-1})=(5,r)(4,r^2)=x^5n\equiv x^5$. So
$K/\langle n\rangle$, which has order $16$ because $K=\{(i,r^js^k):i\equiv j
\pmod 2\}$ has order $32$, is generated by two elements satisfying the defining
relations of $M_{16}$, whence $K/\langle n\rangle\cong M_{16}$ and $M_{16}\in HS(\mathbb Z_8\times
D_8)\subseteq V\vee V'$. The group $M_{16}$ is subdirectly irreducible, its
unique minimal normal subgroup being $\langle a^4\rangle$, and it is in neither
$V$ (non-abelian) nor $V'$ (exponent $8$, while $V'$ has exponent $4$). The
verbal subgroups are $\theta_V(M_{16})=M_{16}'M_{16}^8=\langle a^4\rangle$ and
$\theta_{V'}(M_{16})\supseteq M_{16}^4=\langle a^4\rangle$, with equality
because $M_{16}/\langle a^4\rangle\cong\mathbb Z_4\times\mathbb Z_2\in V'$.
Thus condition (i) fails for $A=M_{16}$, and by Theorem~\ref{thm:B}
the quotient map $M_{16}\to\mathbb Z_4\times\mathbb Z_2$ lies in
$W_V\cap W_{V'}$ but not in $W_{V\vee V'}$. Groups are congruence permutable,
so it is distributivity that is missing.
\end{ex}

\begin{ex}[distributive, not permutable]\label{ex:lattices}
Let $\sV$ be the variety of bounded lattices $(L,\wedge,\vee,0,1)$ with an
additional constant $c$, which is congruence distributive, and let $V$ and $V'$
be the subvarieties defined by $c=0$ and by $c=1$. Let $B=\mathbf 2\times\mathbf 2$
with $c=(0,1)$; it is the product of a $V$-algebra and a $V'$-algebra, so
$B\in V\vee V'$, and $\theta_V(B)=\Cg((0,1),(0,0))$ is the kernel of the first
projection while $\theta_{V'}(B)=\Cg((0,1),(1,1))$ is the kernel of the second.
Let $C=\{(0,0),(0,1),(1,1)\}$ be the three-element chain, a subalgebra of $B$
containing $c$. Then $\theta_V(C)$ identifies $(0,0)$ with $(0,1)$ and
$\theta_{V'}(C)$ identifies $(0,1)$ with $(1,1)$, so the inclusion
$g\colon C\hookrightarrow B$ induces bijections $C/\theta_V(C)\to B/\theta_V(B)$
and $C/\theta_{V'}(C)\to B/\theta_{V'}(B)$, both two-element lattices, while $g$
is not surjective: $C$ is a $V$-dense and $V'$-dense proper subalgebra, (ii)
fails, and $g\in W_V\cap W_{V'}\setminus W_{V\vee V'}$. Condition (i) holds
for all algebras of $V\vee V'$ by distributivity. The verbal congruences of
$C$ do not permute: $(1,1)$ and $(0,0)$ are related by
$\theta_V(C)\vee\theta_{V'}(C)$ but not by $\theta_V(C)\circ\theta_{V'}(C)$,
and $C$ is a proper subalgebra of its pullback, which is $B$.
\end{ex}

\begin{ex}[non-permuting verbal congruences, equality holds]\label{ex:unary}
Let $\sV$ be the variety of mono-unary algebras $(A,f)$ with $f^4=f^2$, and
$V\colon f^2x=x$, $V'\colon f^2x=fx$. The identities $f^px=f^qx$ true in both
are those with $p\equiv q\pmod 2$ and $p,q\ge1$ (or $p=q$), so
$V\vee V'$ is defined by $f^3x=fx$. An algebra $B$ of $V\vee V'$ consists of its
\emph{core} $f(B)$, on which $f$ is an involution, and \emph{hairs}
$x\notin f(B)$ with $f(x)\in f(B)$. The congruence $\theta_V(B)$, generated by
the pairs $(f^2x,x)$, identifies each hair $x$ with the core element $f^2x$ and
nothing else; $\theta_{V'}(B)$, generated by the pairs $(f^2x,fx)$, collapses
each $2$-cycle of the core to a point and nothing else. Hence
$\theta_V(B)\cap\theta_{V'}(B)=\Delta$, which is (i). For (ii), let
$C\subseteq B$ be $V$-dense and $V'$-dense. Density for $V$ means that
$C/\theta_V(C)$, which is the core of $C$, maps onto the core of $B$, so
$C\supseteq f(B)$; density for $V'$ means that the hairs of $C$ are all the
hairs of $B$, since hairs are not identified by $\theta_{V'}$. Thus $C=B$, and
$W_{V\vee V'}=W_V\cap W_{V'}$ by Theorem~\ref{thm:B}. Yet the verbal
congruences do not permute on the three-element algebra $A=\{0,1,2\}$ with
$f(0)=1$, $f(1)=2$, $f(2)=1$: $\theta_V(A)$ identifies $0$ with $2$,
$\theta_{V'}(A)$ identifies $1$ with $2$, and $(1,0)\in\theta_{V'}\circ\theta_V$
but $(1,0)\notin\theta_V\circ\theta_{V'}$; $A$ is a proper subalgebra of its
pullback $A/\theta_V\times A/\theta_{V'}$, which has four elements and in
which $A$ is not $V'$-dense.
\end{ex}

\begin{rem}
Condition (i) for all $A\in V\vee V'$ says that $V\vee V'$ consists of
subdirect products of algebras of $V$ and of $V'$, which Kowalski and Paoli
\cite{KP} study for disjoint subvarieties; in the congruence distributive case
it holds for all pairs by Jónsson's lemma. Condition (ii) is a statement about
dense subalgebras which permutability implies but which, by
Example~\ref{ex:unary}, is weaker. We do not know a characterization of the
pairs $(V,V')$ with $W_{V\vee V'}=W_V\cap W_{V'}$ in terms of the lattice
$\Sub(\sV)$ alone.
\end{rem}

\section{Lattices of logics and of $\lambda$-theories}\label{sec:logic}

For an algebraizable logic $\mathcal L$ with equivalent algebraic semantics the
variety $\sV_{\mathcal L}$, the lattice of axiomatic extensions of $\mathcal L$
is dually isomorphic to $\Sub(\sV_{\mathcal L})$, an extension $L$ corresponding
to the subvariety $\sV_L$ of algebras validating its theorems \cite{BP}. Write
$W_L:=W_{\sV_L}$, the class of homomorphisms of $\sV_{\mathcal L}$-algebras
inverted by the reflection into $\sV_L$; $L\le L'$ gives $W_L\subseteq W_{L'}$.

\begin{cor}\label{cor:logics}
Let $\mathcal L$ be an algebraizable logic whose equivalent algebraic semantics
$\sV_{\mathcal L}$ is congruence distributive and congruence permutable. Then
$L\mapsto W_L$ is a lattice embedding of the lattice of axiomatic extensions of
$\mathcal L$ into $\Same(\sV_{\mathcal L})$: $W_{L\wedge L'}=W_L\cap W_{L'}$ and
$W_{L\vee L'}=W_L\vee W_{L'}$. This applies to
\begin{enumerate}
\item the lattice $\NExt(\mathbf K)$ of normal modal logics, with
$\sV_{\mathcal L}$ the modal algebras \cite{CZ};
\item the lattice of intermediate logics, with Heyting algebras;
\item the lattice of substructural logics over the full Lambek calculus
$\mathbf{FL}$, with $\mathbf{FL}$-algebras, and its intervals of fuzzy logics
over $\mathbf{MTL}$, $\mathbf{BL}$, \L{}ukasiewicz and G\"odel logic, with
$\mathbf{MTL}$-, $\mathbf{BL}$-, $\mathbf{MV}$- and Gödel algebras.
\end{enumerate}
\end{cor}
\begin{proof}
Theorem~\ref{thm:A} and Corollary~\ref{cor:arith} with $\sV=\sV_{\mathcal L}$, using the
dual isomorphism $L\mapsto\sV_L$; the meet of logics is the join of varieties
and conversely. Each of the varieties listed has a lattice reduct, hence is
congruence distributive, and a Mal'cev term: for modal algebras, which have a
Boolean reduct, $x\oplus y\oplus z$ with $\oplus$ symmetric difference; for
Heyting algebras and their expansions the standard term; for
$\mathbf{FL}$-algebras $p(x,y,z)=x/(z\backslash y\wedge 1)\wedge z/(x\backslash
y\wedge 1)$ \cite{JT,BT,GJKO}, which also serves all their subvarieties.
\end{proof}

\begin{rem}
In the framework of \cite{Levelshift}, where $(\Same(C),\subseteq)$ is read as
a Kripke frame whose worlds are the levels of sameness, Corollary~\ref{cor:logics}
says that the lattice of normal modal logics sits inside the frame of modal
algebras as a sublattice: the logics themselves are levels of sameness for the
algebras of the base logic.
\end{rem}

\begin{cor}\label{cor:lambda}
Let $\lT$ be the lattice of $\lambda$-theories and $\mathbf{LAA}$ the variety of
lambda abstraction algebras. The assignment $T\mapsto W_T:=W_{\sV_T}$, where
$\sV_T$ is the subvariety generated by the term algebra $\Lambda/T$, is an
injective map $\lT\to\Same(\mathbf{LAA})$ preserving order and binary joins:
$W_{T\vee T'}=W_T\vee W_{T'}$; and always $W_{T\wedge T'}\subseteq W_T\cap W_{T'}$.
\end{cor}
\begin{proof}
By \cite{Salibra00} the lattice $\lT$ is isomorphic to the lattice of
equational theories of $\mathbf{LAA}$, i.e.\ dually isomorphic to
$\Sub(\mathbf{LAA})$, by $T\mapsto\sV_T$; joins of theories correspond to meets
of varieties, and Theorem~\ref{thm:A} applies.
\end{proof}

\begin{que}\label{que:lambda}
Is $T\mapsto W_T$ a lattice embedding $\lT\hookrightarrow\Same(\mathbf{LAA})$,
i.e.\ does $W_{T\wedge T'}=W_T\cap W_{T'}$ hold for all $\lambda$-theories? By
Theorem~\ref{thm:B} this requires every subdirectly irreducible algebra
of $\sV_T\vee\sV_{T'}$ to satisfy $T$ or $T'$. Corollary~\ref{cor:arith} does not
apply: lambda abstraction algebras satisfy no non-trivial lattice identity in
their congruence lattices \cite{LS04}, and $\lT$ itself is not modular
\cite{Salibra01}.
\end{que}

\begin{figure}[ht]
\[
\begin{tikzcd}[row sep=small]
\NExt(\mathbf K)\arrow[r,hook,"\text{lattice}"] & \Same(\mathbf{MA}) &
\Ext(\mathbf{FL})\arrow[r,hook,"\text{lattice}"] & \Same(\mathbf{FL})\\
\lT\arrow[r,hook,"\vee\text{-semilattice}"'] & \Same(\mathbf{LAA}) &
\mathrm{Top}(S)\arrow[r,hook,"{\text{lattice}, S\text{ finite}}"'] & \Same(\Presh(S))
\end{tikzcd}
\]
\caption{Lattices of theories inside sameness lattices: logics (Corollary~\ref{cor:logics}),
$\lambda$-theories (Corollary~\ref{cor:lambda}), Grothendieck topologies \cite{Finitejoin}.}
\end{figure}

\section*{Funding}
No funding was received for this work.

\section*{Declarations}
\noindent\textbf{CRediT authorship contribution statement.}\quad Ryota Shibusawa:
Conceptualization, Methodology, Formal analysis, Investigation, Software,
Validation, Visualization, Writing -- original draft, Writing -- review \& editing.

\noindent\textbf{Competing interests.}\quad The author declares that he has no
competing interests.

\begin{thebibliography}{99}
\bibitem{BP} W.~J.~Blok, D.~Pigozzi, \emph{Algebraizable logics}, Mem. Amer.
Math. Soc. \textbf{396} (1989).
\bibitem{BT} K.~Blount, C.~Tsinakis, \emph{The structure of residuated
lattices}, Internat. J. Algebra Comput. \textbf{13} (2003), 437--461.
\bibitem{BS} S.~Burris, H.~P.~Sankappanavar, \emph{A Course in Universal
Algebra}, Springer, 1981.
\bibitem{CZ} A.~Chagrov, M.~Zakharyaschev, \emph{Modal Logic}, Oxford University
Press, 1997.
\bibitem{GJKO} N.~Galatos, P.~Jipsen, T.~Kowalski, H.~Ono, \emph{Residuated
Lattices: An Algebraic Glimpse at Substructural Logics}, Elsevier, 2007.
\bibitem{Jonsson} B.~Jónsson, \emph{Algebras whose congruence lattices are
distributive}, Math. Scand. \textbf{21} (1967), 110--121.
\bibitem{JT} P.~Jipsen, C.~Tsinakis, \emph{A survey of residuated lattices},
in: Ordered Algebraic Structures (J.~Martínez, ed.), Kluwer, 2002, 19--56.
\bibitem{KP} T.~Kowalski, F.~Paoli, \emph{Joins and subdirect products of
varieties}, Algebra Universalis \textbf{65} (2011), 371--391.
\bibitem{LS04} S.~Lusin, A.~Salibra, \emph{The lattice of lambda theories},
J. Logic Comput. \textbf{14} (2004), 373--394.
\bibitem{Salibra00} A.~Salibra, \emph{On the algebraic models of lambda
calculus}, Theoret. Comput. Sci. \textbf{249} (2000), 197--240.
\bibitem{Salibra01} A.~Salibra, \emph{Nonmodularity results for lambda
calculus}, Fund. Inform. \textbf{45} (2001), 379--392.
\bibitem{Same} R.~Shibusawa, \emph{The sameness lattice of a category:
weak-equivalence structures, Galois functoriality, and comparison intervals},
preprint (2026), archived at DOI 10.5281/zenodo.22823772.
\bibitem{Levelshift} R.~Shibusawa, \emph{The level-shift logic of a category: modal correspondence for sameness lattices, quotient models with contingent identity, and completeness},
preprint (2026), archived at DOI 10.5281/zenodo.22823775.
\bibitem{Postnikov} R.~Shibusawa, \emph{Postnikov truncations as joins of sheaf
and homotopy equivalences}, preprint (2026), archived at DOI 10.5281/zenodo.22956181.
\bibitem{Topjoin} R.~Shibusawa, \emph{Meets and joins of Grothendieck
topologies in the sameness lattice}, preprint (2026), archived at DOI
10.5281/zenodo.22978446.
\bibitem{Finitejoin} R.~Shibusawa, \emph{Grothendieck topologies on a finite
category form a sublattice of the lattice of two-out-of-three classes},
preprint (2026), archived at DOI 10.5281/zenodo.23204943.
\end{thebibliography}

\end{document}
